\chapter{Explicit-Euler diffusion}
\label{chap:dev:explicit-euler-diffusion}

\begin{quote}
Phase metadata. Spec dir:
\passthrough{\lstinline!specs/2026-05-02-phase-5-explicit-euler/!}.
Implements \passthrough{\lstinline!solve::explicitEuler!} (generic) and
\passthrough{\lstinline!solve::diffuse!} (heat-equation driver with CFL
check). Version target: \passthrough{\lstinline!0.5.0!}.
\end{quote}

\subsection{1. Theory}\label{theory}

Phase 5 introduces \textbf{time integration} to the library. The
continuous problem is the heat equation on a periodic Cartesian domain:

\[
\partial_t \psi(\mathbf{x}, t) \;=\; D\, \nabla^2 \psi(\mathbf{x}, t),
\qquad \psi(\mathbf{x}, 0) = \psi_0(\mathbf{x}),
\]

with \(\psi : \Omega \times [0, T] \to \mathbb{R}\), \(D > 0\), and
periodic boundary conditions on
\(\Omega = [0, L_x) \times [0, L_y) \times [0, L_z)\).

\textbf{Forward (explicit) Euler} discretizes time as
\(\psi^{n+1} = \psi^n + dt \cdot \mathrm{rhs}(\psi^n)\). For the heat
equation, \(\mathrm{rhs}(\psi) = D \nabla^2 \psi\).

The eigenfunctions of the continuous Laplacian on this periodic domain
are tensor-product trig modes (Phase 1 §1). They are also eigenfunctions
of the \textbf{discrete} Laplacian \(\nabla^2_h\), which is convenient:
the continuous solution and the semi-discrete (continuous-time) solution
and the fully-discrete (forward-Euler) solution can all be written
analytically as products of factors of the form \(e^{-\lambda t}\) or
\((1 - dt\,\lambda)^n\). This makes the trig-IC test in §2 a direct
algebraic comparison.

\textbf{Stability} is the new ingredient. The discrete Laplacian is
bounded below (its smallest eigenvalue is the most negative); explicit
Euler is only conditionally stable, and the bound on \(dt\) comes from a
von Neumann analysis of the discrete amplification factor.

\subsection{2. Math derivation}\label{math-derivation}

\subsubsection{Discrete eigenvalues of the Laplacian (recap from Phase
1)}\label{discrete-eigenvalues-of-the-laplacian-recap-from-phase-1}

For \(\psi_{\mathbf{k}} = \prod_a \sin(k_a x_a)\) with integer
wavenumbers \(k_a\), the Phase-1 stencil gives (per axis):

\[
\frac{\psi_{ijk + \hat{a}} - 2 \psi_{ijk} + \psi_{ijk - \hat{a}}}{h_a^2}\Big|_{\psi_{\mathbf{k}}} \;=\; -\frac{2(1 - \cos(k_a h_a))}{h_a^2}\,\psi_{\mathbf{k}}.
\]

Summing over axes,

\[
\nabla^2_h \psi_{\mathbf{k}} \;=\; \lambda_h(\mathbf{k})\,\psi_{\mathbf{k}}, \qquad \lambda_h(\mathbf{k}) \;=\; -\sum_a \frac{2(1 - \cos(k_a h_a))}{h_a^2}.
\]

\subsubsection{Forward-Euler amplification
factor}\label{forward-euler-amplification-factor}

Because \(\psi_{\mathbf{k}}\) is an eigenfunction of \(\nabla^2_h\), the
discrete update on this mode is a scalar:

\[
\psi^{n+1}_{\mathbf{k}} \;=\; \psi^n_{\mathbf{k}} + dt \cdot D \lambda_h(\mathbf{k})\,\psi^n_{\mathbf{k}} \;=\; \bigl(1 + D\,dt\,\lambda_h(\mathbf{k})\bigr)\,\psi^n_{\mathbf{k}}.
\]

The \textbf{amplification factor} is
\(g(\mathbf{k}) = 1 + D\,dt\,\lambda_h(\mathbf{k})\). Since
\(\lambda_h(\mathbf{k}) \le 0\) for all \(\mathbf{k}\) (the Laplacian is
negative semi-definite), \(g \le 1\) automatically. The lower bound
\(g \ge -1\) is the binding stability requirement:

\[
1 + D\,dt\,\lambda_h(\mathbf{k}) \;\ge\; -1
\quad\Longleftrightarrow\quad
D\,dt\,|\lambda_h(\mathbf{k})| \;\le\; 2.
\]

Maximizing over \(\mathbf{k}\) --- each \(1 - \cos(k_a h_a)\) is
maximized at \(2\) (when \(k_a h_a = \pi\)) --- gives the \textbf{CFL
bound}:

\[
\max_{\mathbf{k}} |\lambda_h(\mathbf{k})| \;=\; \sum_a \frac{4}{h_a^2}, \qquad \boxed{\,D\,dt\,\sum_a \frac{1}{h_a^2} \;\le\; \frac{1}{2}.\,}
\]

For isotropic \(h\), this reduces to \(D\,dt/h^2 \le 1/(2d)\) --- the
form quoted in the roadmap. The implementation uses the per-axis sum
because \passthrough{\lstinline!Grid!} allows anisotropic spacing.

\subsubsection{Convergence on the trig
IC}\label{convergence-on-the-trig-ic}

For \(\psi_0 = \sin x \sin y \sin z\) on \([0, 2\pi]^3\) (so
\(\mathbf{k} =
(1, 1, 1)\)), the analytical solution is
\(\psi(t) = e^{-3 D t}\,\psi_0\). After \(n\) explicit-Euler steps the
discrete solution is

\[
\psi^n \;=\; (1 + D\,dt\,\lambda_h)^n\,\psi_0, \qquad \lambda_h \;=\; -\sum_a \frac{2(1 - \cos h_a)}{h_a^2}.
\]

For \(h_a = h\) small, \(\lambda_h = -3 + h^2/4 + O(h^4)\) (using
\(1 - \cos h = h^2/2 - h^4/24 + \ldots\) per axis times 3 axes;
verifying:
\(3 \cdot (h^2/2 - h^4/24) \cdot 2/h^2 = 3 - h^2/4 + O(h^4)\), so
\(\lambda_h = -3 + h^2/4 + O(h^4)\)). Hence
\(1 + D\,dt\,\lambda_h = 1 - 3 D\,dt + (D\,dt)\,h^2/4 + O(dt h^4)\).

Comparing \((1 + D\,dt\,\lambda_h)^n\) to \(e^{-3 D t}\) at
\(t = n\,dt\):

\begin{itemize}
\tightlist
\item
  The \textbf{time-discretization} error from \((1 + x)^n\) vs
  \(e^{n x}\) is \(O(dt)\) globally for explicit Euler.
\item
  The \textbf{spatial} error from \(\lambda_h - (-3) = h^2/4 + O(h^4)\)
  contributes an extra \(O(h^2)\).
\item
  When the CFL is at its limit, \(dt \sim h^2/D\), so the time error
  inherits the same \(O(h^2)\) scaling as the spatial error. The
  combined error is \textbf{\(O(h^2)\)} under fixed-\(T\) refinement.
\end{itemize}

This is the basis for
\passthrough{\lstinline!DiffusionTest.TrigEigenfunctionSecondOrderConvergence!}.
The log-log slope test allows up to \(2.5\) rather than \(2.2\) because
the \(O(dt)\) time-error term contributes a small bonus when \(dt\) is
below the CFL limit (we use \(dt = 0.4 \cdot h^2/(6D)\), well under the
bound).

\subsection{3. Algorithm}\label{algorithm}

\subsubsection{Generic integrator}\label{generic-integrator}

\href{../../../include/gridcalc/solve/ExplicitEuler.hpp}{\passthrough{\lstinline!include/gridcalc/solve/ExplicitEuler.hpp!}}.

\begin{lstlisting}[language={C++}]
template <class Rhs>
inline void explicitEuler(core::Field<double>& psi, Rhs&& rhs, double dt, int n_steps);
\end{lstlisting}

A SFINAE-style \passthrough{\lstinline!static\_assert!} validates that
\passthrough{\lstinline!Rhs!} is callable as
\passthrough{\lstinline!Field<double>(const Field<double>\&)!}. Each
step:

\begin{enumerate}
\def\labelenumi{\arabic{enumi}.}
\tightlist
\item
  \passthrough{\lstinline!Field<double> tmp = rhs(psi)!} --- one fresh
  allocation.
\item
  \passthrough{\lstinline!psi[n] += dt * tmp[n]!} for
  \passthrough{\lstinline!n!} in \passthrough{\lstinline![0, N)!} ---
  raw-pointer loop.
\end{enumerate}

The integrator runs entirely on \passthrough{\lstinline!psi.data()!} and
\passthrough{\lstinline!tmp.data()!}; it deliberately does \textbf{not}
use \passthrough{\lstinline!Field::operator()!} here because the
periodic-wrap policy contributes nothing inside the in-place
accumulation (we touch every element exactly once, in storage order).

\subsubsection{Diffusion driver}\label{diffusion-driver}

\href{../../../include/gridcalc/solve/Diffusion.hpp}{\passthrough{\lstinline!include/gridcalc/solve/Diffusion.hpp!}}.

\begin{lstlisting}[language={C++}]
inline void diffuse(core::Field<double>& psi, double D, double dt, int n_steps);
\end{lstlisting}

Validates \passthrough{\lstinline!D, dt, n\_steps >= 0!}, then calls
\passthrough{\lstinline!checkExplicitDiffusionCFL!} (which throws on
violation), then forwards to
\passthrough{\lstinline!explicitEuler(psi, \&diff::laplacian, D * dt, n\_steps)!}.
The \passthrough{\lstinline!D * dt!} fold is the algebraic identity
\(\psi^{n+1} = \psi^n + dt(D\nabla^2 \psi^n) = \psi^n + (D\,dt)\nabla^2 \psi^n\);
it lets the explicit-Euler step accept an unscaled Laplacian RHS without
a per-step element-wise scaling loop, at the cost of
\passthrough{\lstinline!dt!} losing its physical-time meaning
\textbf{inside \passthrough{\lstinline!explicitEuler!}'s frame}. The
user-visible \passthrough{\lstinline!dt!} in the
\passthrough{\lstinline!diffuse!} API is unaffected.

\subsubsection{Per-step allocation}\label{per-step-allocation}

Each \passthrough{\lstinline!explicitEuler!} step allocates one fresh
\passthrough{\lstinline!Field<double>!} (the RHS return value) and frees
it at the end of the step. For a 100-step run on a \(64^3\) field, this
is 100 alloc/free pairs of \textasciitilde2 MB each. The allocator's
freelist makes this much cheaper than 100 fresh
\passthrough{\lstinline!mmap!}s in practice (most modern
\passthrough{\lstinline!malloc!} implementations recycle the buffer).
Phase 20 will switch to a persistent scratch via a non-allocating RHS
shape (e.g.,
\passthrough{\lstinline!void rhs(const Field\&, Field\& out)!}), which
is the standard approach for tight time-stepping loops.

\subsubsection{CFL check}\label{cfl-check}

\passthrough{\lstinline!checkExplicitDiffusionCFL(grid, D, dt)!}
computes \(\mathrm{cfl} = D\,dt \sum_a (1/h_a^2)\) and throws
\passthrough{\lstinline!std::invalid\_argument!} if
\(\mathrm{cfl} > 0.5\). The message includes the actual values for
diagnostics --- easy to re-compute the right
\passthrough{\lstinline!dt!} from the printed line.

\subsection{4. Design decisions}\label{design-decisions}

\begin{itemize}
\tightlist
\item
  \textbf{In-place mutation} rather than return-fresh. Time integration
  doesn't naturally produce a single output; the standard pattern is to
  evolve a state vector. Returning a fresh
  \passthrough{\lstinline!Field<double>!} per call (let alone per step)
  would double memory churn for no gain. Phase 1--4 operators return
  fresh because they're snapshot-style queries; Phase 5 is the first
  time the iterator pattern dominates.
\item
  \textbf{Throw on CFL violation} rather than warn. Running a 100-step
  explicit Euler past its stability bound produces NaN / Inf within a
  few steps; silent overflow is a debugging trap. A throw with a
  descriptive message lets the caller catch, halve
  \passthrough{\lstinline!dt!}, and retry, which is the expected
  workflow.
\item
  \textbf{No generic integrator interface yet.} Phase 6 introduces
  \passthrough{\lstinline!solve::Integrator!} and
  \passthrough{\lstinline!RK4!}. Anticipating that interface here would
  ship a class hierarchy that gets refactored a phase later. Keep
  \passthrough{\lstinline!explicitEuler!} as a free function; the Phase
  6 refactor will adapt it to consume the new interface.
\item
  \textbf{Trig-IC convergence test rather than periodic-Gaussian.} The
  roadmap's literal ``Gaussian initial condition'' requires summing
  Gaussian image charges over all periodic shifts (a 3D Jacobi theta
  function) for an analytical reference; that's out of proportion to the
  rest of Phase 5's surface area. The trig eigenfunction has a
  closed-form analytical solution that's just \(e^{-3 D t}\) times the
  IC, which makes the comparison trivial to read in the test. The
  Gaussian IC still appears in \passthrough{\lstinline!GaussianSanity!}
  as a no-NaN / mass-conserved / peak-decreased smoke test.
\item
  \textbf{Per-axis CFL form}. Anisotropy-aware. See
  \passthrough{\lstinline!requirements.md!}.
\end{itemize}

\subsection{5. References}\label{references}

{[}1{]} Strikwerda, J. C. (2004). \emph{Finite Difference Schemes and
Partial Differential Equations} (2nd ed.). SIAM. ISBN 978-0-89871-639-9.
\url{https://doi.org/10.1137/1.9780898717938}. §6.1 covers the von
Neumann analysis of forward-Euler discretization of the heat equation,
the derivation reproduced in §2 of this note. §10 gives the
multi-dimensional extension and the per-axis sum form of the CFL bound.

{[}2{]} LeVeque, R. J. (2007). \emph{Finite Difference Methods for
Ordinary and Partial Differential Equations}. SIAM. ISBN
978-0-89871-783-9. \url{https://doi.org/10.1137/1.9780898717839}. §9.2
(forward-Euler heat equation) and §10.2 (von Neumann analysis); cited
from Phase 1 as well but the heat-equation chapter is new context here.

{[}3{]} Wikipedia: \emph{Heat equation} --- \emph{Periodic boundary
conditions and Fourier series solution}.
\url{https://en.wikipedia.org/wiki/Heat_equation\#The_heat_equation_on_the_circle}
(permanent URL). Documents the Fourier-series solution
\(\psi(t) = \sum_{\mathbf{k}} \hat\psi_{\mathbf{k}}(0)\,e^{-D|\mathbf{k}|^2 t}\,e^{i\mathbf{k}\cdot\mathbf{x}}\)
and the eigenfunction decay used in the trig-IC test.

{[}4{]} Press, W. H., Teukolsky, S. A., Vetterling, W. T., \& Flannery,
B. P. (2007). \emph{Numerical Recipes: The Art of Scientific Computing}
(3rd ed.). Cambridge University Press. ISBN 978-0-521-88068-8. §17.2
covers the ``forward-Euler is conditionally stable'' framing and the
practical recommendation to use higher-order time integrators (RK4 in
Phase 6) or implicit methods (Phase 19) for non-trivial diffusion
problems.
